%=========== ΛΥΣΗ - 19ο ΘΕΜΑ Δ ===========
\begin{thema}{Δ}
\begin{erwthma}
\item Η $f$ είναι παραγωγίσιμη στο $\mathbb{R}$ και για κάθε $x\in\mathbb{R}$ είναι
\begin{align*}
f'(x)
&=\frac{2x}{x^2+1}-1\\
&=\frac{2x-x^2-1}{x^2+1}\\
&=-\frac{(x-1)^2}{x^2+1}.
\end{align*}
Επειδή $x^2+1>0$, έχουμε
\[
f'(x)\leq0
\]
για κάθε $x\in\mathbb{R}$, με την ισότητα να ισχύει μόνο για $x=1$. Επομένως η $f$ είναι \textbf{γνησίως φθίνουσα} στο $\mathbb{R}$ και άρα συνάρτηση 1-1.
Πράγματι, είναι γνησίως φθίνουσα στα διαστήματα $(-\infty,1]$ και $[1,+\infty)$, ενώ για $x_1<1<x_2$ ισχύει
\[
f(x_1)>f(1)>f(x_2).
\]

Έχουμε
\begin{align*}
f(e^x-1)=f(0)
&\Leftrightarrow e^x-1=0\\
&\Leftrightarrow e^x=1\\
&\Leftrightarrow {x=0}.
\end{align*}

\item Για κάθε $x\in\mathbb{R}$ είναι
\begin{align*}
f''(x)
&=\left(\frac{2x}{x^2+1}-1\right)'\\
&=\frac{2(x^2+1)-4x^2}{(x^2+1)^2}\\
&=\frac{2(1-x^2)}{(x^2+1)^2}.
\end{align*}
Επειδή $(x^2+1)^2>0$, το πρόσημο της $f''$ είναι το πρόσημο της παράστασης
\[
1-x^2=(1-x)(1+x).
\]
Επομένως:
\begin{itemize}
\item $f''(x)>0$ για $-1<x<1$,
\item $f''(x)=0$ για $x=-1$ ή $x=1$,
\item $f''(x)<0$ για $x<-1$ ή $x>1$.
\end{itemize}
Άρα η $f$ είναι \textbf{κυρτή} στο $[-1,1]$ και \textbf{κοίλη} στα διαστήματα $(-\infty,-1]$ και $[1,+\infty)$.

Επειδή η $f''$ αλλάζει πρόσημο στα $x=-1$ και $x=1$, η $C_f$ έχει σημεία καμπής τα
\[
{K_1(-1,\ln2+1)}
\]
και
\[
{K_2(1,\ln2-1)}.
\]

\item Έχουμε
\[
f(0)=0
\quad\text{και}\quad
f'(0)=-1.
\]
Άρα η εφαπτομένη της $C_f$ στο σημείο $A(0,0)$ έχει εξίσωση
\begin{align*}
y-f(0)&=f'(0)(x-0)\\
\Leftrightarrow {y=-x}.
\end{align*}

Για κάθε $x\in\mathbb{R}$ είναι
\[
f(x)-(-x)=\ln(x^2+1)\geq0,
\]
αφού $x^2+1\geq1$. Επομένως
\[
{f(x)\geq-x}
\]
για κάθε $x\in\mathbb{R}$, άρα η $C_f$ βρίσκεται πάνω από την εφαπτομένη $y=-x$. Η ισότητα ισχύει μόνο για $x=0$.

\item Έχουμε
\begin{align*}
\lim_{x\to0}\frac{f(x)+x}{x^2}
&=\lim_{x\to0}\frac{\ln(x^2+1)}{x^2}\\
&\xlongequal[\text{\eng{DLH}}]{\frac{0}{0}}
\lim_{x\to0}\frac{\frac{2x}{x^2+1}}{2x}\\
&=\lim_{x\to0}\frac{1}{x^2+1}\\
&={1}.
\end{align*}
\end{erwthma}
\end{thema}
%=========== ΤΕΛΟΣ ΛΥΣΗΣ - 19ο ΘΕΜΑ Δ ===========
